\documentclass[11pt]{article}
\usepackage{fontspec}
\setmainfont{Times New Roman}
\setsansfont{Arial}
\setmonofont{Consolas}
\usepackage{amsmath,amssymb,mathtools}
\usepackage{geometry}
\usepackage{microtype}
\geometry{a4paper,margin=2.35cm}
\setlength{\parindent}{1.25em}
\setlength{\parskip}{0pt}
\emergencystretch=25em
\sloppy
\hbadness=10000
\vbadness=10000
\hfuzz=2pt
\vfuzz=10pt
\raggedbottom

\begin{document}

% Унит: Папер27 в001. Дирецт Интерславиц Латин транслатион фром те Герман финал-аудитед соурце слице.
\section*{27. Хилбертове числа в теорији идеалов}

\begin{center}
J. Ber. d. DMV 34 (1925), стр. 101
\end{center}

Е. Noether: Хилбертове числа в теорији идеалов. Под Хилбертовыми числами обче разумејут се така числа, кторе односет се к изчисльеньу остатковых класов идеалов. Сам Hilbert в својој „карактеристичној функцији“ дал функцију, ктора за идеалы полиномној области од некојеј границы далье дава точно число линеарно независных остатковых класов. За нижше степенове числа Macaulay дополнил ту функцију творјачоју функцијеју, ктору Ostrowski могл експлицитно израчунати. Але Macaulay је такоже подал числа другого вида, кторе се показали како најбитнејше за обчу теорију идеалов, ибо можно јих схватити абстрактно. Согласно обшчим разложным теоремам достаточно јест дати дефиницију за примарне идеалы $\mathfrak q$. Иде о число линеарно независных остатковых класов взходно к польу остатковых класов асоциованого простого идеала $\mathfrak p$. Та числа разпадајут се на подсистемы, кторе всакократ сут дане числом неразложимых компонентов идеалов $\mathfrak q,\mathfrak q/\mathfrak p,\mathfrak q/\mathfrak p^2,\ldots$. Число $i$-того подсистема јест такоже карактеризовано числом членов композицијного реда, кторы иде од $\mathfrak q/\mathfrak p^i$ к $\mathfrak q/\mathfrak p^{i-1}$. Целы композицијны ред и јего Хилбертове числа творјат еквивалент представјеньа примарных идеалов како потенциј простых идеалов, кторо в обшчем случају -- напр. в конечных порјадках алгебраичного числового польа -- вече не важи.

\end{document}

